\chapter{Dal problema al programma}
\label{chap:dal-problema}

Il capitolo precedente ha mostrato che cosa sa fare la macchina: eseguire,
velocissima, istruzioni elementari scritte in linguaggio macchina. Questo
capitolo percorre la strada nell'altro senso, partendo da un
\textbf{problema} da risolvere: come si descrive la soluzione (l'algoritmo),
come la si scrive in un linguaggio di alto livello (il programma) e come il
programma diventa qualcosa che la CPU può eseguire.

\section{Programmare}

La programmazione è l'attività che porta a sviluppare un software. Non ogni
procedimento, però, può diventare un programma: deve poter essere
\textbf{descritto in modo non ambiguo}, deve \textbf{raggiungere l'obiettivo
in un tempo finito} e deve essere \textbf{deterministico}. Un procedimento
con queste caratteristiche si chiama \term{algoritmo}: una ricetta di cucina
ben scritta, o il metodo per risolvere un'equazione di secondo grado.

Programmare significa descrivere al calcolatore \textit{come} risolvere un
problema. Il processore mette a disposizione molte operazioni di base
(somma, sottrazione, operazioni logiche, spostamenti di valori, letture e
scritture verso le periferiche), ma, come abbiamo visto, conviene esprimersi
in un linguaggio di alto livello e lasciare a un programma il compito di
tradurre.

\section{I linguaggi di alto livello}

Un linguaggio di alto livello (\textit{HLL}, \textit{High Level Language})
ha tre caratteristiche:

\begin{itemize}
  \item un \textbf{alto livello di astrazione}: le istruzioni hanno un
        significato più complesso di quelle della macchina, e non serve
        sapere come il processore le realizza (\code{A + B} somma due
        valori, punto);
  \item viene \textbf{tradotto} in linguaggio macchina (o in assembly) da
        un programma traduttore, ed è quindi indipendente dal processore e
        dal sistema operativo;
  \item è più \textbf{simile al linguaggio umano}, quindi più facile da
        scrivere e da ricordare (in molti linguaggi \code{print X} mostra il
        valore di \code{X}).
\end{itemize}

Comunque avvenga la traduzione, il programmatore scrive le istruzioni in un
file di testo: il \term{codice sorgente} (o \textit{programma sorgente}, o
semplicemente \textit{programma}). Per tradurlo ci sono due strategie,
l'\textbf{interprete} e il \textbf{compilatore}.

\subsection{Interprete e compilatore}

Un \term{interprete} traduce le istruzioni del codice sorgente \textbf{una
alla volta} in linguaggio macchina, e ciascuna viene eseguita subito dalla
CPU, prima di passare alla successiva. Un \term{compilatore}, invece,
traduce \textbf{tutto} il codice sorgente in linguaggio macchina e lo salva
in un file, il \term{programma eseguibile}; l'esecuzione avviene dopo, ed è
un'operazione separata (figura~\ref{fig:interprete}).

\begin{figure}[H]
  \centering
  \begin{tikzpicture}[font=\Lato\scriptsize,
      src/.style={draw=neutral!55!white, fill=neutral!8!white, minimum width=2.1cm,
        minimum height=0.42cm, font=\Lato\tiny},
      sw/.style={fill=primary!65!white, text=white, minimum width=2.1cm,
        minimum height=0.85cm, font=\Lato\scriptsize\bfseries, align=center},
      ese/.style={draw=neutral!55!white, fill=neutral!14!white, minimum width=2.2cm,
        minimum height=0.85cm, font=\Lato\tiny, align=center}]
    % interprete
    \node[font=\Lato\small\bfseries, anchor=west] at (-1.2,1.1) {Interprete};
    \draw[neutral!40!white, fill=neutral!3!white] (-1.2,0.7) rectangle (1.2,-2.6);
    \node[font=\Lato\tiny\bfseries] at (0,0.5) {FILE SORGENTE};
    \node[src] (i1) at (0,0) {istruzione 1};
    \node[src] (i2) at (0,-0.6) {istruzione 2};
    \node at (0,-1.25) {\vdots};
    \node[src] (in) at (0,-1.9) {istruzione N};
    \node[font=\FiraCodeNL\tiny] at (0,-2.4) {programma.py};
    \node[sw] (s1) at (3.5,0) {interprete};
    \node[sw] (s2) at (3.5,-1.2) {interprete};
    \node[ese] (e1) at (6.3,0) {esecuzione\\istruzione 1};
    \node[ese] (e2) at (6.3,-1.2) {esecuzione\\istruzione 2};
    \node at (3.5,-2.0) {\vdots};
    \draw[wtarrow] (i1.east) -- (s1.west);
    \draw[wtarrow] (i2.east) -- ++(0.6,0) |- (s2.west);
    \draw[wtarrow] (s1) -- (e1);
    \draw[wtarrow] (s2) -- (e2);
    \node[mnota, text width=4.6cm, anchor=north west] at (-1.2,-2.8)
      {la traduzione avviene \textbf{a ogni esecuzione}, un'istruzione alla volta};
    % compilatore
    \begin{scope}[xshift=9.0cm]
      \node[font=\Lato\small\bfseries, anchor=west] at (-1.2,1.1) {Compilatore};
      \draw[neutral!40!white, fill=neutral!3!white] (-1.2,0.7) rectangle (1.2,-2.6);
      \node[font=\Lato\tiny\bfseries] at (0,0.5) {FILE SORGENTE};
      \node[src] at (0,0) {istruzione 1};
      \node[src] at (0,-0.6) {istruzione 2};
      \node at (0,-1.25) {\vdots};
      \node[src] at (0,-1.9) {istruzione N};
      \node[font=\FiraCodeNL\tiny] at (0,-2.4) {programma.c};
      \node[sw] (c) at (2.95,-0.95) {compilatore};
      \node[ese, fill=accent!12!white, draw=accent!50!black, minimum width=1.8cm] (x) at (5.3,-0.95) {file\\eseguibile};
      \node[ese, minimum width=1.8cm] (r) at (7.6,-0.95) {esecuzione};
      \draw[wtarrow] (1.2,-0.95) -- (c);
      \draw[wtarrow] (c) -- (x);
      \draw[wtarrow] (x) -- (r);
      \node[mnota, text width=6.0cm, anchor=north west] at (-1.2,-2.8)
        {la traduzione avviene \textbf{una volta sola}, prima; poi il file
         eseguibile si esegue quante volte si vuole};
    \end{scope}
  \end{tikzpicture}
  \caption{Interprete (a sinistra) e compilatore (a destra).}
  \label{fig:interprete}
\end{figure}

Il C è un linguaggio \textbf{compilato}; Python e JavaScript, per esempio,
sono di solito interpretati. Le due strategie hanno pregi e difetti
opposti:

\begin{center}
\renewcommand{\arraystretch}{1.45}
\begin{tabular}{@{}>{\raggedright\arraybackslash}p{2.8cm}>{\raggedright\arraybackslash}p{6.6cm}>{\raggedright\arraybackslash}p{6.6cm}@{}}
  \toprule
  & \textbf{Programma interpretato} & \textbf{Programma compilato} \\
  \midrule
  \textbf{velocità} & più lento: a ogni esecuzione l'interprete deve ritradurre
    le istruzioni & più veloce: è già tradotto; e il compilatore sa
    \textit{ottimizzare} il codice, per renderlo più veloce o più piccolo \\
  \textbf{chi deve avere che cosa} & l'utente deve avere l'interprete e saperlo
    usare (si può nascondere, ma c'è) & all'utente basta l'eseguibile (un
    doppio clic); solo il programmatore ha bisogno del compilatore \\
  \textbf{codice sorgente} & serve all'utente, quindi è visibile e
    modificabile da chiunque & non serve per eseguire: si può non
    distribuirlo, proteggendo il lavoro (risalire al sorgente è possibile, ma
    difficile) \\
  \bottomrule
\end{tabular}
\end{center}

\subsection{Una via di mezzo: il bytecode}

Alcuni linguaggi, come Java e Python, compilano il sorgente in un file
intermedio, non eseguibile direttamente, scritto in una sorta di assembly
generico chiamato \term{bytecode}, che si traduce molto rapidamente nel
codice macchina di un processore reale. Su ogni computer si installa un
interprete del bytecode (per Java, la \textit{Java Virtual Machine}): così
lo stesso file gira su macchine con linguaggi macchina diversi
(figura~\ref{fig:bytecode}).

\begin{figure}[H]
  \centering
  \begin{tikzpicture}[font=\Lato\scriptsize,
      b/.style={wtnode, font=\Lato\scriptsize, minimum height=0.9cm, text width=2.6cm}]
    \node[b] (s) at (0,0) {sorgente\\{\FiraCodeNL Programma.java}};
    \node[b, draw=accent!60!black, fill=accent!8!white] (c) at (3.6,0) {compilatore\\{\FiraCodeNL javac}};
    \node[b, draw=orange-gradient-6!80!black, fill=orange-gradient-1!30!white] (bc) at (7.2,0) {bytecode\\{\FiraCodeNL Programma.class}};
    \foreach \y/\m in {1.15/{JVM per Windows (x86-64)}, 0/{JVM per macOS (ARM)}, -1.15/{JVM per Linux}} {
      \node[b, draw=green-gradient-6, fill=green-gradient-1!45!white, text width=3.4cm] (j) at (11.6,\y) {\m};
      \draw[wtarrow] (bc.east) -- (j.west);
    }
    \draw[wtarrow] (s) -- (c);
    \draw[wtarrow] (c) -- (bc);
  \end{tikzpicture}
  \caption{Un solo file di bytecode, un interprete diverso per ogni tipo di
           macchina.}
  \label{fig:bytecode}
\end{figure}

\section{Le librerie}

In un linguaggio di alto livello il programmatore non deve programmare da sé
le funzioni di base: calcolare una radice quadrata o una tangente, stampare
un numero sullo schermo, leggere un valore dalla tastiera. Queste operazioni
sono state programmate e compilate una volta per tutte da chi produce il
compilatore, e sono a disposizione sotto forma di \term{funzioni} (per
esempio \code{tan(x)} calcola la tangente di \code{x}, \code{printf} stampa).

Il loro codice, già tradotto in linguaggio macchina, è raccolto in file
chiamati \term{librerie}: collezioni di funzioni pronte all'uso. Un
programma C “vero” è quindi fatto di due parti: il codice scritto da noi
e il codice delle librerie che usa.

\section{Come nasce un eseguibile}

Torniamo al compilatore. Finora lo abbiamo trattato come una scatola
chiusa, che riceve il sorgente e restituisce l'eseguibile; in realtà la
traduzione di un programma C avviene in più passaggi, ciascuno svolto da uno
strumento diverso. Conoscerli aiuta a capire i messaggi di errore, che
cambiano a seconda del passaggio in cui qualcosa è andato storto.

\subsection{Compilazione e linking}

La creazione di un eseguibile (in inglese \textit{build}) avviene in due
fasi (figura~\ref{fig:build}):

\begin{passi}
  \item \textbf{compilazione}: il sorgente viene tradotto in codice
        macchina, ma le funzioni di libreria che usa sono ancora
        \textbf{mancanti}. Il risultato è un file intermedio, il
        \term{file oggetto} (estensione \code{.o}), che non si può
        eseguire;
  \item \textbf{linking} (collegamento): uno o più file oggetto e una o più
        librerie (già in codice macchina) vengono \textbf{uniti} dal
        \term{linker}, che aggiunge le parti mancanti e produce un unico
        file eseguibile.
\end{passi}

\begin{figure}[H]
  \centering
  \begin{tikzpicture}[font=\Lato\scriptsize,
      file/.style={wtnodeplain, minimum width=2.6cm, minimum height=0.8cm},
      comp/.style={wtnode, draw=accent!60!black, fill=accent!8!white,
        font=\Lato\scriptsize\bfseries, minimum width=2.6cm},
      ogg/.style={wtnode, draw=primary!60!black, fill=primary!6!white,
        font=\Lato\scriptsize, minimum width=2.6cm, minimum height=0.8cm}]
    \node[file] (s1) at (0,0) {sorgente {\FiraCodeNL main.c}\\(contiene \code{main})};
    \node[file] (s2) at (3.6,0) {sorgente {\FiraCodeNL funzioni.c}\\(altre funzioni)};
    \node[comp] (c1) at (0,-1.3) {compilatore};
    \node[comp] (c2) at (3.6,-1.3) {compilatore};
    \node[ogg] (o1) at (0,-2.6) {{\FiraCodeNL main.o}\\oggetto, non eseguibile};
    \node[ogg] (o2) at (3.6,-2.6) {{\FiraCodeNL funzioni.o}\\oggetto, non eseguibile};
    \node[ogg, draw=orange-gradient-6!80!black, fill=orange-gradient-1!30!white] (lib) at (7.6,-2.6) {\textbf{libreria}\\funzioni precompilate};
    \node[wtnode, draw=green-gradient-6, fill=green-gradient-1!45!white, minimum width=9cm,
          font=\Lato\scriptsize\bfseries] (lk) at (3.8,-4.0) {linker};
    \node[wtnode, draw=red-gradient-6!70!black, fill=red-gradient-1!25!white,
          font=\Lato\scriptsize, minimum width=2.6cm] (ex) at (3.8,-5.3) {\textbf{eseguibile}\\{\FiraCodeNL programma}};
    \foreach \a/\b in {s1/c1, s2/c2, c1/o1, c2/o2} \draw[wtarrow] (\a) -- (\b);
    \draw[wtarrow, draw=green-gradient-6!70!black] (o1.south) -- (o1.south |- lk.north);
    \draw[wtarrow, draw=green-gradient-6!70!black] (o2.south) -- (o2.south |- lk.north);
    \draw[wtarrow, draw=green-gradient-6!70!black] (lib.south) -- (lib.south |- lk.north);
    \draw[wtarrow, line width=1.3pt] (lk) -- (ex);
  \end{tikzpicture}
  \caption{La compilazione produce i file oggetto, il linking li unisce fra
           loro e con le librerie.}
  \label{fig:build}
\end{figure}

Quando nel capitolo~\ref{chap:ambiente} abbiamo scritto
\code{gcc ciao.c}, gcc ha eseguito entrambe le fasi in un colpo solo. Le si
può però separare, ed è istruttivo farlo con il primo programma di esempio
del corso, \code{area\_cerchio.c} (lo vedremo per intero fra poco).
L'opzione \code{-c} (\textit{compile}) chiede a gcc di fermarsi dopo la
compilazione:

\begin{console}
$ gcc -c area_cerchio.c
$ ls -l area_cerchio.c area_cerchio.o
-rw-r--r-- 1 fabio fabio  255 ott 31 10:09 area_cerchio.c
-rw-r--r-- 1 fabio fabio 1960 ott 31 10:09 area_cerchio.o
$ nm area_cerchio.o
                 U __isoc99_scanf
                 U __stack_chk_fail
0000000000000000 T main
                 U printf
\end{console}

Il comando \code{nm} elenca i nomi che compaiono in un file oggetto. La
\code{T} accanto a \code{main} significa che il suo codice è
\textbf{presente} nel file (\textit{text}, cioè codice); la \code{U} accanto
a \code{printf} e \code{scanf} (che compare con il nome interno
\code{\_\_isoc99\_scanf}) significa \textbf{undefined}: il programma le usa,
ma il loro codice non c'è. Sono proprio le “funzioni di libreria ancora
mancanti”. Il linking le aggiunge:

\begin{console}
$ gcc -o area_cerchio area_cerchio.o
$ ./area_cerchio
Inserisci il raggio del cerchio: <<5>>
L'area del cerchio è: 78.54
\end{console}

Dato un file \code{.o} invece di un \code{.c}, gcc salta la compilazione e
fa solo il linking. Separare le due fasi diventa utile quando un programma è
diviso in più file: se ne cambia uno, basta ricompilare quello e rifare il
linking. Lo vedremo nel capitolo~\ref{chap:make}.

\subsection{La catena completa}

La figura~\ref{fig:catena-completa} mostra tutto il percorso di un programma
C, dall'editor fino alla CPU, con gli strumenti coinvolti.

\begin{figure}[H]
  \centering
  \begin{tikzpicture}[font=\Lato\scriptsize,
      st/.style={wtnode, font=\Lato\scriptsize\bfseries, minimum width=2.5cm, minimum height=0.85cm,
        draw=accent!60!black, fill=accent!8!white},
      fl/.style={wtnodeplain, font=\Lato\scriptsize, minimum width=2.5cm, minimum height=0.85cm}]
    \node[st] (ed) at (0,0) {editor\\{\mdseries (VS Code)}};
    \node[fl] (src) at (3.4,0) {programma sorgente\\{\FiraCodeNL prova.c}};
    \node[st] (pre) at (6.8,0) {preprocessore};
    \node[st] (cmp) at (10.2,0) {compilatore};
    \node[fl] (obj) at (13.6,0) {programma oggetto\\{\FiraCodeNL prova.o}};
    \node[st] (lnk) at (13.6,-2.0) {linker};
    \node[fl, draw=orange-gradient-6!80!black, fill=orange-gradient-1!30!white] (lib) at (13.6,-3.6) {librerie\\{\FiraCodeNL libc, libm\dots}};
    \node[fl] (exe) at (10.2,-2.0) {programma eseguibile\\{\FiraCodeNL prova}};
    \node[st] (ld) at (6.8,-2.0) {loader\\{\mdseries (sistema operativo)}};
    \node[fl, draw=red-gradient-6!70!black, fill=red-gradient-1!20!white] (mem) at (3.4,-2.0) {programma\\in memoria};
    \node[st, draw=primary!65!black, fill=primary!12!white] (cpu) at (0,-2.0) {CPU};
    \node[st, draw=green-gradient-6, fill=green-gradient-1!45!white] (dbg) at (3.4,-3.6) {debugger\\{\mdseries (gdb)}};
    \foreach \a/\b in {ed/src, src/pre, pre/cmp, cmp/obj, obj/lnk, lnk/exe, exe/ld, ld/mem, mem/cpu}
      \draw[wtarrow] (\a) -- (\b);
    \draw[wtarrow] (lib) -- (lnk);
    \draw[wtarrowacc, dashed, <->] (dbg) -- node[wtlab, right=2pt, fill=none] {osserva e controlla} (mem);
  \end{tikzpicture}
  \caption{Dal sorgente al processore: gli strumenti (in viola) e i file che
           producono.}
  \label{fig:catena-completa}
\end{figure}

\begin{itemize}
  \item L'\textbf{editor} serve a scrivere il sorgente.
  \item Il \term{preprocessore} lavora sul testo prima della compilazione
        vera e propria: esegue le righe che cominciano con \code{\#}.
        \code{\#include <stdio.h>} viene sostituita dall'intero contenuto
        del file \code{stdio.h}; \code{\#define PI 3.14159} fa sì che ogni
        \code{PI} del programma diventi \code{3.14159}.
  \item Il \textbf{compilatore} traduce in codice macchina e produce il
        file oggetto; il \textbf{linker} lo unisce alle librerie.
  \item Il \term{loader} è la parte del sistema operativo che, quando
        scriviamo \code{./prova}, copia l'eseguibile in memoria e fa partire
        la CPU dalla sua prima istruzione (capitolo~\ref{chap:calcolatore}).
  \item Il \term{debugger} permette di eseguire il programma un passo alla
        volta, fermarlo e guardare il contenuto delle variabili: è lo
        strumento per la \textit{messa a punto}.
\end{itemize}

Il lavoro del preprocessore si può vedere con l'opzione \code{-E}, che
ferma gcc subito dopo questa fase e stampa il testo risultante. Il nostro
programma di 15 righe diventa di oltre 800, quasi tutte provenienti da
\code{stdio.h}; in fondo c'è il nostro codice, con \code{PI} già sostituito:

\begin{console}
$ gcc -E area_cerchio.c | wc -l
828
$ gcc -E area_cerchio.c | tail -6
    scanf("%f", &raggio);

    area = 3.14159 * raggio * raggio;

    printf("L'area del cerchio è: %.2f\n", area);
\end{console}

(Il simbolo \code{|} è una \textit{pipe}: passa l'output di un comando a un
altro. \code{wc -l} conta le righe, \code{tail -6} mostra le ultime 6.)

\subsection{Librerie statiche e dinamiche}

Il linking può collegare una libreria in due modi.

Con una \term{libreria statica}, il codice delle funzioni usate viene
\textbf{copiato} dentro il file eseguibile. Le stesse funzioni possono
essere usate da più programmi, ma ognuno ne contiene una copia
(figura~\ref{fig:librerie}, a sinistra).

Con una \term{libreria dinamica}, nell'eseguibile \textbf{non} viene inserito
il codice delle funzioni, ma solo un \textbf{riferimento}: il nome del file
della libreria e il nome della funzione. Su Windows sono le \textit{DLL}
(\textit{Dynamic Link Libraries}), su Linux e macOS le \textit{shared
libraries} (file \code{.so} e \code{.dylib}). Quando parte il primo
programma che usa una libreria dinamica, il sistema operativo la carica in
memoria; se poi parte un altro programma che usa la stessa libreria, la
trova già in memoria e la \textbf{condivide} (figura~\ref{fig:librerie}, a
destra).

\begin{figure}[H]
  \centering
  \begin{tikzpicture}[font=\Lato\scriptsize,
      app/.style={draw=primary!65!black, fill=primary!10!white, line width=0.7pt,
        minimum width=1.6cm, minimum height=0.55cm, font=\Lato\tiny\bfseries},
      lib/.style={draw=orange-gradient-6!80!black, fill=orange-gradient-1!35!white,
        line width=0.7pt, minimum width=1.6cm, minimum height=1.3cm, font=\Lato\tiny,
        align=center}]
    % statiche
    \node[font=\Lato\scriptsize\bfseries] at (1.85,1.2) {librerie statiche};
    \node[app] (a1) at (0.9,0.45) {applicazione A};
    \node[lib, anchor=north] at ([yshift=0.7pt]a1.south) {copia della\\libreria};
    \node[app] (a2) at (2.8,0.45) {applicazione B};
    \node[lib, anchor=north] at ([yshift=0.7pt]a2.south) {copia della\\libreria};
    % dinamiche
    \node[font=\Lato\scriptsize\bfseries] at (8.1,1.2) {librerie dinamiche};
    \node[app] (b1) at (6.9,0.45) {applicazione A};
    \node[draw=primary!65!black, fill=white, line width=0.7pt, minimum width=1.6cm,
          minimum height=0.45cm, anchor=north, font=\Lato\tiny] (r1) at ([yshift=0.7pt]b1.south) {riferimento};
    \node[app] (b2) at (9.3,0.45) {applicazione B};
    \node[draw=primary!65!black, fill=white, line width=0.7pt, minimum width=1.6cm,
          minimum height=0.45cm, anchor=north, font=\Lato\tiny] (r2) at ([yshift=0.7pt]b2.south) {riferimento};
    \node[lib] (sh) at (8.1,-1.75) {libreria\\condivisa\\(una sola,\\in memoria)};
    \draw[mrif] (r1.south) -- (sh.north west);
    \draw[mrif] (r2.south) -- (sh.north east);
  \end{tikzpicture}
  \caption{A sinistra ogni eseguibile contiene la sua copia della libreria;
           a destra gli eseguibili contengono solo un riferimento a un'unica
           copia, condivisa.}
  \label{fig:librerie}
\end{figure}

I vantaggi delle librerie dinamiche sono evidenti: gli eseguibili sono più
piccoli, un programma parte più in fretta se la libreria è già in memoria e,
se una funzione va aggiornata (una nuova versione, la correzione di un
errore), basta sostituire la libreria, senza ritoccare i programmi. Il
rovescio della medaglia: se la libreria manca o è danneggiata, il programma
non funziona.

gcc usa le librerie dinamiche, se non si chiede diversamente. Il comando
\code{ldd} mostra da quali dipende un eseguibile: il nostro dipende da
\code{libc.so.6}, la libreria standard del C. Con l'opzione \code{-static}
si ottiene invece un collegamento statico, e la differenza di dimensione
parla da sola:

\begin{console}
$ ldd area_cerchio
	linux-vdso.so.1 (0x00007f1b90918000)
	libc.so.6 => /lib/x86_64-linux-gnu/libc.so.6 (0x00007f1b90600000)
	/lib64/ld-linux-x86-64.so.2 (0x00007f1b9091a000)
$ gcc -static -o area_statico area_cerchio.o
$ ls -l area_cerchio area_statico
-rwxr-xr-x 1 fabio fabio  16072 ott 31 10:09 area_cerchio
-rwxr-xr-x 1 fabio fabio 915784 ott 31 10:09 area_statico
\end{console}

Lo stesso programma pesa 16 kilobyte con la libreria dinamica e quasi un
megabyte con quella statica, che contiene dentro di sé tutto il codice di
\code{printf}, \code{scanf} e delle funzioni che queste usano. (Su macOS
\code{-static} non è disponibile per i programmi normali, e al posto di
\code{ldd} si usa \code{otool -L}.)

\section{Gli algoritmi}

\subsection{Che cos'è un algoritmo}

Possiamo ora precisare la definizione del capitolo precedente. Un
algoritmo è la \textbf{descrizione operativa} di come risolvere un problema:
un \textbf{elenco finito di istruzioni elementari}, da eseguire in
successione. Le singole istruzioni devono:

\begin{itemize}
  \item richiedere un \textbf{tempo finito};
  \item \textbf{non essere ambigue};
  \item avere un \textbf{effetto osservabile}.
\end{itemize}

Inoltre l'algoritmo deve produrre un \textbf{risultato}, sempre lo stesso a
partire dalle stesse condizioni iniziali, e le risorse che richiede devono
essere \textbf{ragionevoli} per la tecnologia attuale: il tempo complessivo
di esecuzione, la memoria massima necessaria, eventuali dispositivi
periferici particolari.

Il termine nasce in informatica, ma si applica a qualunque procedimento.
Cambiano le istruzioni elementari, non le caratteristiche:

\begin{center}
\renewcommand{\arraystretch}{1.3}
\begin{tabular}{@{}ll@{}}
  \toprule
  \textbf{Algoritmo per\dots} & \textbf{un'istruzione elementare} \\
  \midrule
  risolvere un'equazione di secondo grado & calcola $b^2 - 4ac$ \\
  preparare una ricetta                   & sbatti 4 albumi \\
  montare un armadio IKEA                 & unisci i pannelli 12 e 13 con una vite C \\
  avviare il motore di un aereo           & inserisci i magneti \\
  installare un'app sul cellulare         & tocca il pulsante Installa \\
  \bottomrule
\end{tabular}
\end{center}

Sono invece istruzioni \textbf{non adeguate}:

\begin{itemize}
  \item “calcola tutte le cifre di $\pi$”: le cifre sono infinite, quindi
        l'istruzione non termina mai (tempo non finito);
  \item “aggiungi alcuni degli ingredienti”: quali? quanti? (ambigua).
\end{itemize}

E un algoritmo che con gli stessi dati iniziali produce risultati diversi
\textbf{non è corretto}, perché non è deterministico: un'equazione di
secondo grado deve dare sempre le stesse soluzioni per gli stessi valori di
$a$, $b$ e $c$.

\subsection{Quando le risorse non bastano}

Un algoritmo che richiede un tempo irragionevole per la tecnologia attuale,
considerati costi e benefici, \textbf{non risolve il problema}, anche se è
corretto. L'esempio delle lezioni: aprire un file zip protetto da password
senza conoscerla. Un algoritmo c'è, ed è semplicissimo: provare tutte le
combinazioni di caratteri finché non si trova quella giusta (il cosiddetto
attacco a \textit{forza bruta}).

Supponiamo che la password sia di 8 caratteri, scelti fra i 96 caratteri
stampabili del codice ASCII standard. Per il primo carattere ci sono 96
possibilità; per ciascuna, 96 per il secondo, e così via: in tutto
$96^8 \approx 7{,}2 \cdot 10^{15}$ password possibili, cioè circa 7,2
milioni di miliardi. Con un computer che ne prova 34 milioni al secondo,
senza mai fermarsi:

\[
  \frac{7{,}2 \cdot 10^{15}}{3{,}4 \cdot 10^{7}\ \text{al secondo}}
  \approx 2{,}1 \cdot 10^{8}\ \text{secondi} \approx 6{,}7\ \text{anni}.
\]

\warning{Un conto da rifare}{
  Le slide parlano di “quasi 5 anni” per provare tutte le password; il
  calcolo esatto dà circa \textbf{6,7 anni}. In media, però, la password
  giusta si trova dopo averne provate la metà, cioè in circa 3,4 anni. La
  conclusione non cambia: per chi ha bisogno del file oggi, l'algoritmo è
  inutile.}

La cosa più istruttiva è che cosa succede allungando la password: ogni
carattere in più \textbf{moltiplica per 96} il numero di tentativi.

\begin{center}
\renewcommand{\arraystretch}{1.3}
\begin{tabular}{@{}crr@{}}
  \toprule
  \textbf{lunghezza} & \textbf{password possibili} & \textbf{tempo per provarle tutte} \\
  \midrule
  8  & $7{,}2 \cdot 10^{15}$ & circa 6,7 anni \\
  9  & $6{,}9 \cdot 10^{17}$ & circa 645 anni \\
  10 & $6{,}6 \cdot 10^{19}$ & circa 62\,000 anni \\
  \bottomrule
\end{tabular}
\end{center}

Una crescita di questo tipo si chiama \textit{esponenziale}, ed è il motivo
per cui le password lunghe sono sicure: non serve un algoritmo più furbo per
violarle, servirebbe un tempo che nessuno ha.

\section{Come si descrive un algoritmo}

Prima di scrivere un programma conviene descrivere l'algoritmo in una forma
intermedia, più libera di un linguaggio di programmazione. Le forme più
usate sono:

\begin{itemize}
  \item la \textbf{descrizione testuale}, in italiano;
  \item lo \term{pseudo-codice}: una descrizione testuale con un certo
        livello di formalismo, con parole chiave fisse (\code{LEGGI},
        \code{STAMPA}, \code{SE}\dots) e una struttura simile a quella di un
        programma, ma senza le regole rigide di un linguaggio vero;
  \item il \term{diagramma di flusso} (\textit{flow-chart}): una
        rappresentazione grafica, in cui ogni forma ha un significato
        (figura~\ref{fig:forme});
  \item il \textbf{linguaggio di alto livello}, cioè il programma vero e
        proprio.
\end{itemize}

\begin{figure}[H]
  \centering
  \begin{tikzpicture}[font=\Lato\scriptsize]
    \node[fcterm] (a) at (0,0) {Inizio};
    \node[fcio] (b) at (3.4,0) {leggi \code{x}};
    \node[fcproc] (c) at (6.8,0) {\code{y = 2 * x}};
    \node[fcdec] (d) at (10.2,0) {\code{y > 10}?};
    \node[fcsub] (e) at (13.6,0) {ordina};
    \foreach \nodo/\testo in {a/{inizio e fine}, b/{input e output\\(leggere, stampare)},
                       c/{un'istruzione\\(un calcolo)}, d/{una decisione:\\due uscite, sì e no},
                       e/{un sottoprogramma\\(cap.~\ref{chap:strutturata})}}
      \node[mnota, align=center, anchor=north] at ([yshift=-0.3cm]\nodo.center |- d.south) {\testo};
  \end{tikzpicture}
  \caption{Le forme dei diagrammi di flusso.}
  \label{fig:forme}
\end{figure}

Conviene anche organizzare la descrizione in tre parti: l'\textbf{input}
(i dati di partenza), l'\textbf{elaborazione} (i calcoli) e
l'\textbf{output} (i risultati). Vediamo i tre esempi delle lezioni, in
tutte le forme. Il codice C è commentato solo a grandi linee: ogni
dettaglio sarà spiegato nel capitolo~\ref{chap:primi-programmi}. Qui conta
vedere come ogni riga di pseudo-codice diventa una riga di programma.

\subsection{Esempio 1: l'area di un cerchio}

\begin{itemize}
  \item \textbf{Input}: il raggio del cerchio $r$, un numero reale positivo.
  \item \textbf{Elaborazione}: l'area si calcola con la formula
        $A = \pi \cdot r^2$, dove $\pi$ è la costante che rappresenta il
        rapporto fra la circonferenza di un cerchio e il suo diametro.
  \item \textbf{Output}: il valore numerico dell'area.
\end{itemize}

In forma testuale: \textit{acquisisci dall'utente il valore del raggio;
calcola l'area con la formula $\pi \cdot r \cdot r$; visualizza il valore
ottenuto.} In pseudo-codice e come diagramma di flusso:

\begin{affianco}[0.42][c]
\begin{lstlisting}[style=pseudo, name=pseudocodice]
BEGIN
  LEGGI raggio
  AREA = PI * raggio * raggio
  STAMPA AREA
END
\end{lstlisting}
\risultato
\begin{immagine}
\begin{tikzpicture}[font=\Lato\scriptsize]
  \node[fcterm] (a) at (0,0) {Inizio};
  \node[fcio] (b) at (0,-1.0) {leggi \code{raggio}};
  \node[fcproc] (c) at (0,-2.0) {\code{area = PI * raggio * raggio}};
  \node[fcio] (d) at (0,-3.0) {stampa \code{area}};
  \node[fcterm] (e) at (0,-4.0) {Fine};
  \foreach \x/\y in {a/b, b/c, c/d, d/e} \draw[fcfreccia] (\x) -- (\y);
\end{tikzpicture}
\end{immagine}
\end{affianco}

Ed ecco il programma C, con l'esecuzione accanto:

\begin{affianco}[0.55]
\begin{lstlisting}[style=c, name=area_cerchio.c]
#include <stdio.h>
#define PI 3.14159

int main() {
    float raggio, area;

    printf("Inserisci il raggio del cerchio: ");
    scanf("%f", &raggio);

    area = PI * raggio * raggio;

    printf("L'area del cerchio è: %.2f\n", area);

    return 0;
}
\end{lstlisting}
\risultato
\begin{console}
$ gcc -Wall -o area_cerchio area_cerchio.c
$ ./area_cerchio
Inserisci il raggio del cerchio: <<5>>
L'area del cerchio è: 78.54
\end{console}
\end{affianco}

La corrispondenza con lo pseudo-codice è diretta:

\begin{center}
\renewcommand{\arraystretch}{1.35}
\begin{tabular}{@{}>{\FiraCodeNL\small}l>{\FiraCodeNL\small}l>{\raggedright\arraybackslash}p{6.4cm}@{}}
  \toprule
  \normalfont\textbf{pseudo-codice} & \normalfont\textbf{C} & \textbf{nota} \\
  \midrule
  & float raggio, area; & in C bisogna dichiarare le variabili prima di
                          usarle, dicendone il tipo (\code{float}: numero
                          reale) \\
  LEGGI raggio & scanf("\%f", \&raggio); & prima si stampa un messaggio
                          che dice all'utente che cosa scrivere \\
  AREA = PI * raggio * raggio & area = PI * raggio * raggio; & identico;
                          \code{PI} è definito dalla riga \code{\#define} \\
  STAMPA AREA & printf("\dots \%.2f\textbackslash n", area); & \code{\%.2f}
                          significa: un numero reale con 2 decimali \\
  \bottomrule
\end{tabular}
\end{center}

\subsection{Esempio 2: da Fahrenheit a Celsius}

\begin{itemize}
  \item \textbf{Input}: una temperatura in gradi Fahrenheit (un numero
        reale).
  \item \textbf{Elaborazione}: la formula di conversione è
        $\mathit{Celsius} = (\mathit{Fahrenheit} - 32) \cdot 5 / 9$.
  \item \textbf{Output}: la temperatura in gradi Celsius.
\end{itemize}

Lo pseudo-codice ha la stessa forma del primo esempio, e così il programma:

\begin{affianco}[0.5][c]
\begin{lstlisting}[style=pseudo, name=pseudocodice]
BEGIN
  LEGGI fahrenheit
  CELSIUS = (fahrenheit - 32) * 5 / 9
  STAMPA CELSIUS
END
\end{lstlisting}
\risultato
\wtnota{LEGGI diventa \code{scanf}, il calcolo resta identico, STAMPA
diventa \code{printf}.}
\end{affianco}

\begin{affianco}[0.55]
\begin{lstlisting}[style=c, name=conversione.c]
#include <stdio.h>

int main() {
    float fahrenheit, celsius;

    printf("Inserisci la temperatura in Fahrenheit: ");
    scanf("%f", &fahrenheit);

    celsius = (fahrenheit - 32) * 5 / 9;

    printf("La temperatura in Celsius è: %.2f\n", celsius);

    return 0;
}
\end{lstlisting}
\risultato
\begin{console}
$ gcc -Wall -o conversione conversione.c
$ ./conversione
Inserisci la temperatura in Fahrenheit: <<98.6>>
La temperatura in Celsius è: 37.00
\end{console}
\end{affianco}

\warning{L'ordine delle operazioni conta}{
  Matematicamente $(F - 32) \cdot 5 / 9$ e $5 / 9 \cdot (F - 32)$ sono la
  stessa cosa. In C no: scrivendo \code{celsius = 5 / 9 * (fahrenheit - 32);}
  il risultato è sempre \textbf{0}. Il motivo è che \code{5} e \code{9} sono
  numeri interi, e in C la divisione fra interi dà un intero, troncando i
  decimali: \code{5 / 9} vale 0, e 0 per qualunque cosa fa 0. Nella versione
  del programma, invece, \code{(fahrenheit - 32)} è un \code{float}, quindi
  anche le operazioni successive avvengono fra reali. È una delle trappole
  più comuni del C, e la ritroveremo nel capitolo~\ref{chap:primi-programmi}.}

\subsection{Esempio 3: il segno di un numero}

Il terzo esempio introduce una \textbf{scelta}: il programma deve fare cose
diverse a seconda del dato.

\begin{itemize}
  \item \textbf{Input}: un numero intero.
  \item \textbf{Elaborazione}: se il numero è maggiore di zero è positivo;
        se è minore di zero è negativo; altrimenti è nullo.
  \item \textbf{Output}: un messaggio che indica se il numero è positivo,
        negativo o nullo.
\end{itemize}

\begin{affianco}[0.46][c]
\begin{lstlisting}[style=pseudo, name=pseudocodice]
BEGIN
  LEGGI numero
  SE numero > 0 ALLORA
    STAMPA "Il numero è positivo"
  ALTRIMENTI SE numero < 0 ALLORA
    STAMPA "Il numero è negativo"
  ALTRIMENTI
    STAMPA "Il numero è nullo"
  END SE
END
\end{lstlisting}
\risultato
\begin{immagine}
\begin{tikzpicture}[font=\Lato\scriptsize]
  \node[fcterm] (a) at (0,0) {Inizio};
  \node[fcio] (b) at (0,-0.95) {leggi \code{numero}};
  \node[fcdec] (c) at (0,-2.15) {\code{numero > 0}?};
  \node[fcio] (p) at (3.1,-2.15) {stampa “positivo”};
  \node[fcdec] (d) at (0,-3.75) {\code{numero < 0}?};
  \node[fcio] (n) at (3.1,-3.75) {stampa “negativo”};
  \node[fcio] (z) at (0,-5.05) {stampa “nullo”};
  \node[fcterm] (e) at (1.55,-6.05) {Fine};
  \draw[fcfreccia] (a) -- (b);
  \draw[fcfreccia] (b) -- (c);
  \draw[fcfreccia] (c) -- node[fcsi, above] {sì} (p);
  \draw[fcfreccia] (c) -- node[fcno, right] {no} (d);
  \draw[fcfreccia] (d) -- node[fcsi, above] {sì} (n);
  \draw[fcfreccia] (d) -- node[fcno, right] {no} (z);
  \draw[fcfreccia] (p.east) -- ++(0.35,0) |- (e.east);
  \draw[fcfreccia] (n.east) -- ++(0.35,0) |- (e.east);
  \draw[fcfreccia] (z.south) |- (e.west);
\end{tikzpicture}
\end{immagine}
\end{affianco}

Nel diagramma di flusso si vede bene che, qualunque sia il numero, viene
eseguita \textbf{una sola} delle tre stampe: le strade si separano ai rombi
e si ricongiungono prima della fine. In C la scelta si scrive con
\code{if} (“se”) ed \code{else} (“altrimenti”):

\begin{affianco}[0.55]
\begin{lstlisting}[style=c, name=positivo.c]
#include <stdio.h>

int main() {
    int numero;

    printf("Inserisci un numero: ");
    scanf("%d", &numero);

    if (numero > 0) {
        printf("Il numero è positivo.\n");
    } else if (numero < 0) {
        printf("Il numero è negativo.\n");
    } else {
        printf("Il numero è zero.\n");
    }

    return 0;
}
\end{lstlisting}
\risultato
\begin{console}
$ ./positivo
Inserisci un numero: <<7>>
Il numero è positivo.
$ ./positivo
Inserisci un numero: <<-3>>
Il numero è negativo.
$ ./positivo
Inserisci un numero: <<0>>
Il numero è zero.
\end{console}
\wtnota{Per verificare un programma con una scelta, lo si prova con un
dato per ogni strada possibile.}
\end{affianco}

Si noti l'ultima nota: un programma con tre strade va provato \textbf{tre
volte}, una per strada. Un errore nascosto in un ramo che non si prova mai
resta lì, in attesa del giorno sbagliato.

\insintesi{
  Un algoritmo è un elenco finito di istruzioni elementari, non ambigue,
  ciascuna eseguibile in un tempo finito, che produce sempre lo stesso
  risultato dagli stessi dati con risorse ragionevoli. Lo si descrive in
  italiano, in pseudo-codice, con un diagramma di flusso e infine in un
  linguaggio di alto livello. Il programma sorgente va tradotto: un
  \textbf{interprete} lo fa un'istruzione alla volta a ogni esecuzione, un
  \textbf{compilatore} una volta sola. Per il C la traduzione passa per il
  preprocessore, il compilatore (che produce i file oggetto) e il linker
  (che li unisce alle librerie, statiche o dinamiche); al momento
  dell'esecuzione il loader carica il programma in memoria.}

\section{Riferimenti}

\begin{itemize}
  \item \textbf{Deitel \& Deitel}, \textit{Il linguaggio C}: capitoli 1 e 3
        (ambiente di sviluppo, algoritmi, pseudo-codice).
  \item \textbf{GCC --- Overall Options}: le opzioni \code{-E}, \code{-S},
        \code{-c} e \code{-o}. \\
        \urlx{https://gcc.gnu.org/onlinedocs/gcc/Overall-Options.html}
\end{itemize}
